\documentclass[a4paper,12pt]{article}
\usepackage{fontspec}
\usepackage{polyglossia}
\setotherlanguage{english}
\setmainfont{LiberationSans} % or another font that supports Greek
\usepackage{amsmath, amssymb}
\usepackage[most]{tcolorbox}
\usepackage{xcolor}
\usepackage{colortbl}
\usepackage{enumitem}
\usepackage{multirow}
\usepackage[a4paper, margin=1.8cm]{geometry}
\usepackage{fancyhdr}
\usepackage{hyperref}
\usepackage{booktabs}
\usepackage{tikz}
\usepackage{caption}
\pagestyle{fancy}
\renewcommand{\headrulewidth}{0pt}
\fancyhf{} 
\fancyfoot[L]{\footnotesize Theolaos}
\fancyfoot[R]{\small \thepage}

\usetikzlibrary{decorations.pathreplacing}
\usetikzlibrary{arrows.meta}

\setlist[itemize]{left=1.5em, label=\textbullet}
\setlength{\parskip}{0pt}
\setlength{\parindent}{0pt}

\begin{document}

\begin{center}
    \LARGE \textbf{Σημειώσεις Ηλεκτρονικής Φυσικής}\\[0.5ex]
    \large Τυπολόγιο \\
    \normalsize Βρες και άλλες σημειώσεις: \url{https://github.com/theolaos/simeiwseis}
\end{center}

\section*{Νόμος Coulomb}
\vspace{-.4cm}
\hrule
\vspace{0.3cm}

\begin{equation*}
    F = k \frac{Q_1 Q_2}{r^2} \qquad \text{ή σε διανυσματική μορφή} \qquad \vec{F}_{12} = k \frac{Q_1 Q_2}{r^2} \hat{r}
\end{equation*}

\vspace{0.3cm}

% Σχήμα με TikZ
\noindent
\begin{minipage}[c]{0.55\textwidth}
\textbf{Επεξήγηση Μεγεθών:}
\begin{itemize}\small
    \item \textbf{$r$}: Απόσταση των κέντρων των φορτίων $Q_1, Q_2$.
    \item \textbf{$\hat{r}$}: Μοναδιαίο διάνυσμα (από $Q_1 \rightarrow Q_2$).
    \item \textbf{$\vec{F}_{12}$}: Δύναμη από το $Q_1$ στο $Q_2$.
    \item \textbf{$k$}: Ηλεκτρική σταθερά ($k = \frac{1}{4\pi\varepsilon_0}$).
\end{itemize}
\end{minipage}%
\hfill
\begin{minipage}[c]{0.45\textwidth}
\centering
\vspace{0.2cm}
\begin{tikzpicture}[scale=0.85, >=Stealth]
    % Φορτίο Q1 (+)
    \shade[ball color=blue!40] (0,0) circle (0.3cm);
    \node at (0,0) [white] {\textbf{\small $+$}};
    \node[below=5pt] at (0,0) {$Q_1$};

    % Φορτίο Q2 (+)
    \shade[ball color=blue!40] (3.5,0) circle (0.3cm);
    \node at (3.5,0) [white] {\textbf{\small $+$}};
    \node[below=5pt] at (3.5,0) {$Q_2$};

    % Απόσταση r
    \draw[<->, thick, dashed] (0,0.6) -- (3.5,0.6) node[midway, above] {\small $r$};

    % Δυνάμεις (Απωστικές)
    \draw[->, ultra thick, blue!70!black] (3.75,0) -- (4.8,0) node[right] {\small $\vec{F}_{12}$};
    \draw[->, ultra thick, blue!70!black] (-0.25,0) -- (-1.3,0) node[left] {\small $\vec{F}_{21}$};

    % Μοναδιαίο διάνυσμα r_hat
    \draw[->, thick, green!50!black] (0.4, 0) -- (1.3, 0) node[above] {\small $\hat{r}$};
\end{tikzpicture}

\vspace{0.1cm}

% Σχήμα 2: Ετερόσημα Φορτία (Ελκτικές Δυνάμεις)
\begin{tikzpicture}[scale=0.85, >=Stealth]
    % Φορτίο Q1 (+)
    \shade[ball color=blue!40] (0,0) circle (0.3cm);
    \node at (0,0) [white] {\textbf{\small $+$}};
    \node[below=5pt] at (0,0) {$Q_1$};

    % Φορτίο Q2 (-)
    \shade[ball color=red!50] (3.5,0) circle (0.3cm);
    \node at (3.5,0) [white] {\textbf{\small $-$}};
    \node[below=5pt] at (3.5,0) {$Q_2$};

    % Δυνάμεις (Ελκτικές)
    \draw[->, ultra thick, red!70!black] (3.25,0) -- (2.2,0) node[above] {\small $\vec{F}_{12}$};
    \draw[->, ultra thick, blue!70!black] (0.25,0) -- (1.3,0) node[above] {\small $\vec{F}_{21}$};
\end{tikzpicture}
\end{minipage}

\section*{Ηλεκτρικό Πεδίο}
\vspace{-.4cm}
\hrule
\vspace{0.3cm}

\begin{equation*}
    \vec{E} = \frac{\vec{F}}{q} \qquad \text{αλλιώς μόνο το μέτρο} \qquad E = k\frac{Q}{r^2} \qquad \text{ως πρως $F$:} \qquad  \vec{F} = q\vec{E}
\end{equation*}
\noindent
\begin{minipage}[c]{0.42\textwidth}
\textbf{Επεξήγηση Μεγεθών:}
\begin{itemize}\small
    \item \textbf{$\vec{E}$}: Ένταση ηλεκτρικού πεδίου.
    \item \textbf{$Q$}: Πηγή του πεδίου (φορτίο).
    \item \textbf{$q$}: Φορτίο δοκιμής.
    \item \textbf{$\vec{F}$}: Ηλεκτρική δύναμη στο $q$.
    \item \textbf{$r$}: Απόσταση από το φορτίο $Q$.
\end{itemize}
\end{minipage}%
\hfill
\begin{minipage}[c]{0.55\textwidth}
\centering
% Σχήματα Ηλεκτρικού Πεδίου
\begin{tikzpicture}[scale=0.85, >=Stealth]
    % --- Θετικό Φορτίο (+) ---
    \begin{scope}[shift={(0,0)}]
        % Ακτίνες πεδίου (προς τα έξω)
        \foreach \angle in {0, 45, 90, 135, 180, 225, 270, 315} {
            \draw[->, thick, blue!70!black] (0,0) -- (\angle:1.2);
        }
        % Φορτίο Q (+)
        \shade[ball color=blue!40] (0,0) circle (0.35cm);
        \node at (0,0) [white] {\textbf{\small $+$}};
        \node[below=30pt] at (0,0) {\small $Q > 0$};
    \end{scope}

    % --- Αρνητικό Φορτίο (-) ---
    \begin{scope}[shift={(3.5,0)}]
        % Ακτίνες πεδίου (προς τα μέσα)
        \foreach \angle in {0, 45, 90, 135, 180, 225, 270, 315} {
            \draw[->, thick, red!70!black] (\angle:1.3) -- (\angle:0.4);
        }
        % Φορτίο Q (-)
        \shade[ball color=red!50] (0,0) circle (0.35cm);
        \node at (0,0) [white] {\textbf{\small $-$}};
        \node[below=30pt] at (0,0) {\small $Q < 0$};
    \end{scope}
\end{tikzpicture}
\end{minipage}



\begin{equation*}
    \text{Αρχή Επαλληλίας:} \ E = E_1 + E_2 + E_3 + \dots + E_n \qquad \text{όπου: $E$ πεδίο, $n$ κάποιου φορτίου} 
\end{equation*}

\vspace{-0.3cm}

\begin{equation*}
    \text{Για συνεχή κατανομή:} \ \vec{E} = \int d\vec{E} \qquad \text{όπου:} \ d\vec{E} = \frac{1}{4\pi\epsilon_0}\cdot\frac{dQ}{r^2}\hat{r} 
\end{equation*}

\section*{Πυκνότητα Φορτίου}
\vspace{-.4cm}
\hrule
\vspace{0.3cm}

\subsection*{Ομοιόμορφη Κατανομή}

\begin{itemize}
    \item Χωρική (όγκος) Πυκνώτητα φορτίου: $\rho = \frac{Q}{V}$ όπου είναι σε μονάδες: $\frac{C}{m^3}$
    \item Επιφανειακή (εμβαδόν) Πυκνώτητα φορτίου: $\sigma = \frac{Q}{A}$ όπου είναι σε μονάδες: $\frac{C}{m^2}$
    \item Γραμμική (ευθεία) Πυκνώτητα φορτίου: $\lambda = \frac{Q}{l}$ όπου είναι σε μονάδες: $\frac{C}{m}$
\end{itemize}

\subsection*{Μη-Ομοιόμορφη Κατανομή}
\begin{itemize}
    \item Στοιχειώδη Χωρική (όγκος) Πυκνώτητα φορτίου: $dQ = \rho dV$
    \item Στοιχειώδη Επιφανειακή (εμβαδόν) Πυκνώτητα φορτίου: $dQ = \sigma dA$
    \item Στοιχειώδη Γραμμική (ευθεία) Πυκνώτητα φορτίου: $dQ = \lambda dl$
\end{itemize}

\section*{Ηλεκτριοκό Πεδίο σε ομοιόμορφη κατανομή φορτίου}
\vspace{-.4cm}
\hrule
\vspace{0.3cm}

Έστω απόσταση από δίσκο $z$ και ακτίνα δίσκου $R$, άμα $z<<R$ τότε ισχύει ότι το Ηλ. Πεδίο είναι:
$$E = \frac{\sigma}{2\epsilon_0}$$

Η συμπεριφορά του ηλεκτρικου πεδίου ως προς την απόσταση συγκριτικά με: (και αντί για δίσκο θα μπορούσες να πεις επιφάνεια)
\begin{itemize}
    \item Σημειακοκύ φορτίου: $E \approx \frac{1}{r^2}$
    \item Ομοιόμορφη γραμμική κατανομή μεγάλου μήκους: $E \approx \frac{1}{r}$
    \item Ομοιόμορφη κοτανομή σε μεγάλο επίπεδο τότε το $E$ είναι ανεξάρτητο της απόστασης.
\end{itemize}

\section*{Ηλεκτροστατική Ισορροπία}
\vspace{-.4cm}
\hrule
\vspace{0.3cm}

Όταν αγωγός τοποθετηθεί σε ηλεκτρικό πεδίο,
μετακινήσεις ηλεκτρονίων στην επιφάνεια του
αγωγού πολώνουν τον αγωγό. Εντός του
αγωγού το πεδίο είναι μηδέν.

Η ανακατανομή γίνεται μέσα σε $10^{-16}$ και μπορεί να
θεωρηθεί ακαριαία.

\section*{Κίνηση Φορτισμένου Σωματιδίου σε Ηλεκτρικό Πεδίο}
\vspace{-.4cm}
\hrule
\vspace{0.3cm}

Να θυμάσε πως $\vec{F} = q\vec{E} = m\vec{a}$, που είναι ο δεύετερος νόμος του νεύτωνα. 
Και να συμβουλευτείς του τύπους της μηχανικής, και την ελεύθερης βολής για τον υπολογισμό
της ταχύτητας του ηλεκτρονίου.

\end{document}